\clearpage
\VerdierPhysicalPageAnchor{0246}%
\begingroup\small\setlength{\parskip}{0pt}\setlength{\emergencystretch}{1.5em}
\setlength{\abovedisplayskip}{4pt plus 2pt minus 2pt}%
\setlength{\belowdisplayskip}{4pt plus 2pt minus 2pt}%

\VerdierSemanticAnchor{ch4d-4-6-4}
\noindent\textbf{4.6.4.}\quad
Let $X^{\bullet\bullet}$ be a double complex of $\mathcal A$
(Chapter~I, 2.1.2). The double complex $X^{\bullet\bullet}$ defines an
object of $\operatorname{comp}(\operatorname{comp}(\mathcal A))$:

\VerdierSemanticAnchor{ch4d-4-6-4-double-complex}
\[
X^{\bullet\bullet}\colon\quad
\cdots\longrightarrow X^{\bullet,n-1}
\xrightarrow{d_2^{\bullet,n-1}}X^{\bullet,n}
\xrightarrow{d_2^{\bullet,n}}X^{\bullet,n+1}
\longrightarrow\cdots.
\]
(We use the notation of (4.1.9). Note the sign!)

\VerdierSemanticAnchor{ch4d-4-6-4-b1}
Denote by $B_1^p(X^{\bullet\bullet})$ the simple complex

\VerdierSemanticAnchor{ch4d-4-6-4-1}
\[
\left\{
\begin{aligned}
B_1^p(X^{\bullet\bullet})^i
  &=\operatorname{Im}(d_1^{p-1,i}),\\
d_{B_1^p(X^{\bullet\bullet})}^{i}
  &\text{ induced by }d_2^{p,i}.
\end{aligned}
\right.
\tag{4.6.4.1}
\]

\VerdierSemanticAnchor{ch4d-4-6-4-z1}
Denote by $Z_1^p(X^{\bullet\bullet})$ the simple complex

\VerdierSemanticAnchor{ch4d-4-6-4-2}
\[
\left\{
\begin{aligned}
Z_1^p(X^{\bullet\bullet})^i
  &=\operatorname{Ker}(d_1^{p,i}),\\
d_{Z_1^p(X^{\bullet\bullet})}^{i}
  &\text{ induced by }d_2^{p,i}.
\end{aligned}
\right.
\tag{4.6.4.2}
\]

\VerdierSemanticAnchor{ch4d-4-6-4-h1}
Denote by $H_1^p(X^{\bullet\bullet})$ the simple complex

\VerdierSemanticAnchor{ch4d-4-6-4-3}
\[
\left\{
\begin{aligned}
H_1^p(X^{\bullet\bullet})^i
  &=Z_1^p(X^{\bullet\bullet})^i/B_1^p(X^{\bullet\bullet})^i,\\
d_{H_1^p(X^{\bullet\bullet})}^{i}
  &\text{ induced by }d_2^{p,i}.
\end{aligned}
\right.
\tag{4.6.4.3}
\]

\VerdierSemanticAnchor{ch4d-4-6-4-simple-as-double}
Finally, if $Y$ is a simple complex of $\mathcal A$, we also denote by
$Y$ the double complex whose components in nonzero second degree are zero
and whose components and first differentials in second degree zero are
those of $Y$.

\VerdierSemanticAnchor{ch4d-4-6-5}
\medskip\noindent\textbf{Definition 4.6.5.}\quad
Let $Y$ be a simple complex of $\mathcal A$. A \emph{Cartan--Eilenberg
injective resolution} of $Y$ is a double complex

\VerdierSemanticAnchor{ch4d-4-6-5-double-complex}
\[
I^{\bullet\bullet}\colon\quad
\cdots\longrightarrow0\longrightarrow I^{\bullet,0}
\longrightarrow I^{\bullet,1}\longrightarrow I^{\bullet,2}
\longrightarrow\cdots,
\]
equipped with a morphism called an \emph{augmentation}

\VerdierSemanticAnchor{ch4d-4-6-5-augmentation}
\[
\varepsilon\colon Y\longrightarrow I^{\bullet\bullet}
\]
such that, for every integer $p$, the simple complexes
$B_1^p(I^{\bullet\bullet})$ (4.6.4.1) and
$H_1^p(I^{\bullet\bullet})$ (4.6.4.3), equipped with the augmentations

\VerdierSemanticAnchor{ch4d-4-6-5-augmentations}
\[
B_1^p(\varepsilon)\colon
\operatorname{Im}(d_Y^{p-1})\longrightarrow B_1^p(I^{\bullet\bullet}),
\]
\[
H_1^p(\varepsilon)\colon
H^p(Y)\longrightarrow H_1^p(I^{\bullet\bullet})
\]
\endgroup
